\section{Design}
\label{sec:design}

\subsection{Components}

HyperPod has three parts: a Go controller on the machine the user sits at, a
Rust agent on every laptop that contributes compute, and a React dashboard.
Figure~\ref{fig:arch} shows how they relate to the \texttt{llama.cpp} processes
that perform the actual inference.

\begin{figure}[htbp]
\centering
\begin{tikzpicture}[font=\scriptsize,
  box/.style={draw=gray!55,rounded corners=2pt,align=center,inner sep=4pt,line width=0.4pt},
  mach/.style={box,text width=3.05cm,align=center},
  rstep/.style={box,text width=2.75cm,align=center},
  lnk/.style={-{Stealth[length=4pt]},gray!65,line width=0.5pt}]

% ---- row 1: dashboard and control plane ----
\node[box,text width=2.1cm] (ui) at (-4.85,3.25)
  {\textbf{Dashboard}\\[1pt]{\tiny telemetry}\\[-1pt]{\tiny placement}};
\node[box,fill=serA!10,draw=serA!65,text width=7.4cm] (cp) at (1.55,3.25)
  {\textbf{Go controller}\\[2.5pt]
   {\tiny registry \ \textbar\ \ measured-VRAM scheduler \ \textbar\ \ supervisor \ \textbar\ \ retry}};
\draw[lnk] (ui) -- node[above,font=\tiny,inner sep=2pt]{HTTP} (cp);

% ---- row 2: machines ----
\node[mach,fill=serA!7,draw=serA!50]        (rog) at (-3.35,1.05)
  {\textbf{ROG} · Linux\\[1pt]{\tiny controller + \code{llama-server}}\\[-1pt]{\tiny RTX 3050 · local CUDA}};
\node[mach,fill=green!6,draw=green!40!black!45] (tuf) at (0.45,1.05)
  {\textbf{TUF} · Windows\\[1pt]{\tiny agent + \code{ggml-rpc-server}}\\[-1pt]{\tiny RTX 4050 · remote CUDA}};
\node[mach,fill=orange!8,draw=orange!60!black!40] (mac) at (4.25,1.05)
  {\textbf{M4 Mac} · macOS\\[1pt]{\tiny agent + \code{ggml-rpc-server}}\\[-1pt]{\tiny remote Metal}};

\draw[lnk] (cp.south -| rog.north) -- node[left,font=\tiny,inner sep=1.5pt]{local} (rog.north);
\draw[lnk] (cp.south -| tuf.north) -- node[right,font=\tiny,inner sep=1.5pt]{RPC / Wi-Fi} (tuf.north);
\draw[lnk] (cp.south -| mac.north) -- node[right,font=\tiny,inner sep=1.5pt]{RPC / Wi-Fi} (mac.north);

% discovery / heartbeat path back to the registry
\draw[lnk,serB,dashed] (rog.south) -- ++(0,-0.45) -- (6.35,-0.45) -- (6.35,3.25) -- (cp.east);
\node[font=\tiny,text=serB,anchor=south] at (1.0,-0.42) {discovery · heartbeat · telemetry};

% ---- row 3: recovery sequence ----
\node[rstep,fill=serB!7,draw=serB!50]            (f1) at (-3.35,-1.75)
  {accelerator stops\\[1pt]{\tiny \code{llama-server} aborts}};
\node[rstep,fill=serA!8,draw=serA!50]            (f2) at (0.45,-1.75)
  {endpoint excluded\\[1pt]{\tiny rebuild on survivors}};
\node[rstep,fill=green!7,draw=green!40!black!45] (f3) at (4.25,-1.75)
  {same request retried\\[1pt]{\tiny response returned}};
\draw[lnk,serB,dashed] (f1) -- (f2);
\draw[lnk,serA] (f2) -- (f3);
\end{tikzpicture}
\caption{HyperPod components and the recovery path. \texttt{llama.cpp} performs
all inference; the Go controller decides placement, observes it, and replaces
the model process when a worker stops answering. The bottom row is the sequence
measured in Section~\ref{sec:recovery}.}
\label{fig:arch}
\end{figure}

\subsection{Telemetry that refuses to guess}

Each agent samples its own hardware every five seconds and reports the snapshot
with registration, with every heartbeat, and in its discovery broadcast. The
schema is versioned, and the controller rejects a version it does not
understand.

The design rule that mattered most is that every value which can be unavailable
is nullable, and an unavailable value is transmitted as \code{null} rather than
as zero or as a plausible substitute. Three cases from our own hardware show why:

\begin{itemize}[leftmargin=1.4em,itemsep=2pt]
  \item The Windows RTX~4050 reports no GPU power limit; \code{nvidia-smi}
        returns \code{[N/A]}. That field stays \code{null}.
  \item Apple Silicon has one unified memory pool, so there is no separate
        free-video-memory quantity. Reporting system RAM in that field would
        imply a dedicated figure that does not exist, so it stays \code{null}.
  \item macOS exposes no direct equivalent of \code{MemAvailable}. We derive
        available memory from \code{vm\_stat} page classes and name that source
        explicitly, so a reader can judge the estimate.
\end{itemize}

Sources are carried in the payload (\code{nvidia-smi}, \code{windows-cim},
\code{macos-sysctl+vm\_stat}, and so on) and shown in the dashboard, so no
number in the interface is anonymous.

\subsection{Deriving the split from measurement}

When every selected device reports free video memory, the controller computes

\[
w_i \;=\; \mathrm{free\_vram}_i \;-\; R, \qquad R = \SI{512}{MiB},
\]

where $R$ is a fixed reserve held back on each device for runtime buffers, and
the resulting $w_i$ are passed to \code{--tensor-split} in RPC-first order.
Endpoints are sorted before the weights are built, so the same measurements
always produce the same argument string. $R$ is a policy choice, not a
measurement, and it is labelled as such in the interface.

If any required value is missing --- as with the Mac --- the controller supplies
no weights at all and records the strategy as
\code{llama.cpp-auto (remote free VRAM unavailable)}. \texttt{llama.cpp} then
probes the devices itself. We preferred this to deriving a weight from a GPU
name or from total memory, because a wrong weight produces a silently
unbalanced model rather than a visible gap.

\subsection{Choosing the failure signal}

Detecting worker loss was the hardest design decision, and our first two
attempts were wrong.

Port probing was too eager: a TCP connection can succeed against a socket whose
server is no longer serving, and it can fail transiently on a busy Wi-Fi link.
Health polling was no better, because a loaded model can be slow to answer
while remaining healthy.

The signal we settled on is the exit of the supervised \code{llama-server}
process. When an RPC peer is lost, \texttt{llama.cpp} aborts, so the process
exit is both reliable and specific. Two refinements were needed:

\begin{itemize}[leftmargin=1.4em,itemsep=2pt]
  \item \textbf{Loading is not failure.} A model still loading answers health
        checks with 503. An early version treated that as death and killed its
        own load. Readiness is now watched by a separate observer with no
        timeout, and a request that arrives during loading waits and is then
        told the model is still loading.
  \item \textbf{The abort and the client error race.} The HTTP request fails
        milliseconds before the process exit is observed. A failed request
        therefore waits up to five seconds to see whether the process exits
        before it is classified as an ordinary error.
\end{itemize}

On rebuild, every candidate endpoint is probed directly. A worker whose
heartbeat is healthy but whose accelerator is unreachable is excluded, which is
exactly the condition tested in Section~\ref{sec:recovery}.
